\section{Method}
\label{sec:method}

\subsection{Problem Setup}
\label{sec:setup}

Let $\mathcal{D} = \mathcal{D}_{\mathrm{real}} \cup \mathcal{D}_{\mathrm{synth}}$ denote a pretraining corpus containing both human-authored and machine-generated text.
The contamination ratio $\rho = |\mathcal{D}_{\mathrm{synth}}| / |\mathcal{D}|$ is unknown---a realistic assumption for web-scale corpora where synthetic content enters from diverse, uncoordinated sources~\citep{cheng_2025_contamination_survey}.

A pre-trained detector assigns each sample $x$ a synthetic probability $q(x) \in [0,1]$.
Following \citet{drayson_2025_detection_collapse}, we weight each sample as
\begin{equation}
  w(x) = (1 - q(x))^{b},
  \label{eq:weight}
\end{equation}
where $b \geq 0$ controls downweighting intensity: $b = 0$ recovers uniform training, while larger $b$ progressively suppresses high-$q$ samples.
\citet{drayson_2025_detection_collapse} select $b$ via validation-set grid search, obtaining different optimal values for different model scales ($b{=}10$ for ${\leq}360$M parameters, $b{=}1$ for 1.7B).
We instead derive $b$ per-batch from the batch's own $q$-score distribution, requiring no held-out validation data.

\subsection{Dual-Constraint Framework}
\label{sec:dual}

Given a batch $\mathcal{B} = \{x_1, \ldots, x_B\}$ with detector scores $q_i = q(x_i)$ and weights $w_i(b) = (1 - q_i)^{b}$, we define two batch-level diagnostics that jointly determine $b$.

\paragraph{Effective real proportion.}
Let $r_i = \mathbf{1}[q_i \leq 0.5]$ estimate whether sample $x_i$ is human-authored.
The \emph{effective real proportion}
\begin{equation}
  \alpha_{\mathrm{eff}}(b) = \frac{\sum_{i=1}^{B} w_i(b) \cdot r_i}{\sum_{i=1}^{B} w_i(b)}
  \label{eq:alpha}
\end{equation}
measures the share of total weight allocated to likely-real samples.
At $b = 0$, $\alpha_{\mathrm{eff}}$ equals the raw real fraction in the batch; increasing $b$ shifts weight toward low-$q$ samples, raising $\alpha_{\mathrm{eff}}$.

\paragraph{Effective sample size.}
Kish's effective sample size~\citep{vehtari_2024_psis}
\begin{equation}
  \mathrm{ESS}(b) = \frac{\bigl(\sum_{i=1}^{B} w_i(b)\bigr)^{2}}{\sum_{i=1}^{B} w_i(b)^{2}}
  \label{eq:ess}
\end{equation}
quantifies how many equally-weighted samples the weighted batch is equivalent to.
$\mathrm{ESS}(0) = B$; as $b$ grows, a few high-$(1{-}q)$ samples dominate and $\mathrm{ESS}$ drops toward~1.

\paragraph{Monotonicity.}
The two diagnostics exhibit opposing monotonicity in $b$:

\begin{proposition}
\label{prop:monotone}
For any batch containing both real and synthetic samples with distinct $q$-values, $\alpha_{\mathrm{eff}}(b)$ is strictly increasing in $b$.
\end{proposition}

\begin{proof}[Proof sketch]
Let $\pi_i(b) = w_i / \sum_j w_j$ denote normalized weights.
Differentiating $\alpha_{\mathrm{eff}} = \sum_i \pi_i r_i$ yields
$\frac{d\alpha_{\mathrm{eff}}}{db} = \mathrm{Cov}_{\pi}\!\bigl(r_i,\, \ln(1{-}q_i)\bigr)$.
Real samples ($r_i = 1$) have $q_i \leq 0.5$ and thus higher $\ln(1{-}q_i)$, so this covariance is non-negative for any detector with better-than-random discrimination.
Full derivation in Appendix~\ref{app:proofs}.
\end{proof}

\noindent $\mathrm{ESS}(b)$ is monotonically decreasing in $b$---larger exponents concentrate the weight distribution, a standard property of exponential tilting~\citep{kimura_2024_iw_survey}.

\paragraph{Dual constraints.}
We impose two requirements on $b$:
\begin{enumerate}[nosep,leftmargin=*]
  \item \textbf{Quality floor}: $\alpha_{\mathrm{eff}}(b) \geq \tau$.
  Convergence under contaminated training improves with the effective real fraction~\citep{wang_2026_survive_contamination}; $\tau$ sets the minimum acceptable quality level.
  By Proposition~\ref{prop:monotone}, this constraint defines a \emph{lower bound} $b_\tau = \inf\{b \geq 0 : \alpha_{\mathrm{eff}}(b) \geq \tau\}$.
  \item \textbf{Stability ceiling}: $\mathrm{ESS}(b) \geq n_{\min}$.
  Overly concentrated weights yield high-variance gradients; $n_{\min}$ prevents the batch from degenerating into a single effective sample.
  Since $\mathrm{ESS}$ is decreasing, this defines an \emph{upper bound} $b_{\mathrm{ESS}} = \sup\{b \geq 0 : \mathrm{ESS}(b) \geq n_{\min}\}$.
\end{enumerate}
The optimal exponent is
\begin{equation}
  b^{*} = \inf\bigl\{b \geq 0 : \alpha_{\mathrm{eff}}(b) \geq \tau\bigr\} \;\;\text{s.t.}\;\; \mathrm{ESS}(b) \geq n_{\min}.
  \label{eq:bstar}
\end{equation}
This is the \emph{minimum intervention principle}: apply the smallest $b$ that achieves the quality floor.
Borderline samples---those with moderate $q$-values---still carry useful training signal; over-downweighting discards this signal unnecessarily.
When both constraints are satisfiable ($b_\tau \leq b_{\mathrm{ESS}}$), $b^{*} = b_\tau$ and $\alpha_{\mathrm{eff}}$ is the binding constraint---$b^{*}$ lies well below the ESS ceiling.
When they conflict ($b_\tau > b_{\mathrm{ESS}}$), the stability ceiling takes precedence: $b^{*} = b_{\mathrm{ESS}}$, accepting a lower $\alpha_{\mathrm{eff}}$ to preserve gradient stability.
Figure~\ref{fig:dual_constraint} illustrates the feasible region and the three regimes.

\subsection{Algorithm}
\label{sec:algorithm}

Algorithm~\ref{alg:bstar} computes $b^{*}$ for each training batch via bisection on the monotone functions $\alpha_{\mathrm{eff}}$ and $\mathrm{ESS}$.
Each batch falls into one of three regimes:
\begin{enumerate}[nosep,leftmargin=*]
  \item \textbf{Clean} ($\alpha_{\mathrm{eff}}(0) \geq \tau$): the quality floor holds at $b = 0$; no reweighting needed.
  \item \textbf{$\alpha_{\mathrm{eff}}$-binding} ($b_\tau \leq b_{\mathrm{ESS}}$): the quality constraint determines $b^{*} = b_\tau$; $\mathrm{ESS}$ remains above $n_{\min}$ with margin.
  \item \textbf{ESS-binding} ($b_\tau > b_{\mathrm{ESS}}$): the stability ceiling caps $b^{*} = b_{\mathrm{ESS}}$.
\end{enumerate}
Across 26{,}458 micro-batches in our main experiment ($\rho{=}0.4$, $\tau{=}0.8$, $n_{\min}{=}2$), 49.9\% are clean, 20.2\% are $\alpha_{\mathrm{eff}}$-binding, and 30.0\% are ESS-binding.
In $\alpha_{\mathrm{eff}}$-binding batches, the solver selects $b^{*} = 1.67$ on average, while an ESS-only strategy would choose $b = 4.15$ on the same batches---2.5$\times$ more aggressive.
The residual $\mathrm{ESS} = 2.60$ in these batches (vs.\ $n_{\min} = 2.0$) confirms that the dual constraint prevents unnecessary over-downweighting.

\begin{algorithm}[t]
\caption{Per-batch adaptive $b^{*}$ selection}
\label{alg:bstar}
\begin{algorithmic}[1]
\Require Batch scores $\{q_i\}_{i=1}^{B}$, indicators $\{r_i\}$, quality threshold~$\tau$, ESS floor~$n_{\min}$, search bound~$b_{\max}$, tolerance~$\varepsilon$
\Ensure Optimal exponent $b^{*}$
\If{$\alpha_{\mathrm{eff}}(0) \geq \tau$}
  \State \Return $0$ \Comment{clean batch}
\EndIf
\State $b_\tau \gets \Call{Bisect}{\alpha_{\mathrm{eff}}(\cdot){=}\tau,\; [0, b_{\max}],\; \varepsilon}$
\If{$\mathrm{ESS}(b_\tau) \geq n_{\min}$}
  \State \Return $b_\tau$ \Comment{$\alpha_{\mathrm{eff}}$-binding}
\Else
  \State $b_{\mathrm{ESS}} \gets \Call{Bisect}{\mathrm{ESS}(\cdot){=}n_{\min},\; [0, b_{\max}],\; \varepsilon}$
  \State \Return $b_{\mathrm{ESS}}$ \Comment{ESS-binding}
\EndIf
\end{algorithmic}
\end{algorithm}

\paragraph{Implementation.}
Given $b^{*}$, we compute normalized weights $\pi_i = w_i(b^{*}) / \sum_j w_j(b^{*})$ and form the weighted loss $\mathcal{L} = \sum_i \pi_i \ell_i$, where $\ell_i$ is the per-sample cross-entropy.
This \emph{loss reweighting} maintains continuous gradients, unlike the SIR resampling of \citet{drayson_2025_detection_collapse} which introduces discrete sampling noise by drawing a subset proportional to $w_i$.
Each bisection requires $\mathcal{O}(B \log(b_{\max}/\varepsilon))$ operations---with $B{=}4$, $b_{\max}{=}20$, and $\varepsilon{=}0.01$, roughly 44 floating-point operations per batch, negligible beside the forward and backward passes.

\subsection{Hyperparameters}
\label{sec:hyperparameters}

The framework introduces two hyperparameters, both with principled defaults set once across all experiments.

\paragraph{Quality threshold $\tau$.}
Convergence theory for contaminated training shows that convergence rate improves with the effective real fraction~\citep{wang_2026_survive_contamination}, motivating a lower bound on $\alpha_{\mathrm{eff}}$.
The specific value of $\tau$ is set empirically.
We find a wide safe range: $\tau \in [0.5, 0.8]$ consistently outperforms uniform training (Section~\ref{sec:results}), with $\tau = 0.8$ achieving the best perplexity.
Above 0.8, $b^{*}$ inflates rapidly (mean $b^{*} = 3.08$ at $\tau = 0.95$), causing over-downweighting.
We use $\tau = 0.8$ as the default throughout.

\paragraph{ESS floor $n_{\min}$.}
We set $n_{\min} = 2$ for batch size $B = 4$.
The constraint $n_{\min} < B$ is necessary: $n_{\min} \geq B$ forces $b^{*} = 0$ on every batch, recovering uniform training.
At the other extreme, $n_{\min} = 1$ effectively removes the stability constraint---mean $b^{*}$ rises to 7.20, with 46.7\% of batches having $b^{*} > 2$, and perplexity degrades by $+0.14$ relative to $n_{\min} = 2$.
In practice, $n_{\min} \in \{2, 3\}$ yields stable results across seeds (Section~\ref{sec:results}).
